% --- Archon physics formalization source begin ---
% archon:physics
% archon:covers PhyXMiniProblems/problem_phyx_mini_0519.lean
% archon:source-report reports/phyx_mini/problem_phyx_mini_0519.source.json

\chapter{Physics problem phyx\_mini\_0519}
\label{ch:PhyXMiniProblems_problem_phyx_mini_0519}

\paragraph{Problem source.}
\#\# Physical scenario

In a set of experiments on a hypothetical one-electron atom, you measure the wavelengths of the photons emitted from transitions ending in the ground state \textbackslash{}( n = 1 \textbackslash{}), as shown in the energy-level diagram in figure. You also observe that it takes \textbackslash{}( 17.50 \textbackslash{}text\{ eV\} \textbackslash{}) to ionize this atom.

\#\# Question

If an electron made a transition from the \textbackslash{}( n = 4 \textbackslash{}) to the \textbackslash{}( n = 2 \textbackslash{}) level, what wavelength of light would it emit?

\#\# Answer choices

- A: A: \textbackslash{}( 323 \textbackslash{}text\{ nm\} \textbackslash{})

- B: B: \textbackslash{}( 456 \textbackslash{}text\{ nm\} \textbackslash{})

- C: C: \textbackslash{}( 423 \textbackslash{}text\{ nm\} \textbackslash{})

- D: D: \textbackslash{}( 378 \textbackslash{}text\{ nm\} \textbackslash{})

\#\# Figure caption (auxiliary; use the image as primary evidence)

The image depicts an energy level diagram for an atom with electron transitions between levels. There are five horizontal lines labeled with quantum numbers \textbackslash{}( n = 1 \textbackslash{}) to \textbackslash{}( n = 5 \textbackslash{}). Each line represents an energy level where an electron can reside.

The transitions are shown as vertical arrows pointing downward, indicating electron transitions from higher energy levels to a lower energy level (n=1). The wavelengths (\textbackslash{}(\textbackslash{}lambda\textbackslash{})) of the emitted photons during these transitions are labeled along the arrows:

- From \textbackslash{}( n = 5 \textbackslash{}) to \textbackslash{}( n = 1 \textbackslash{}), \textbackslash{}(\textbackslash{}lambda = 73.86 \textbackslash{}, \textbackslash{}text\{nm\}\textbackslash{}).
- From \textbackslash{}( n = 4 \textbackslash{}) to \textbackslash{}( n = 1 \textbackslash{}), \textbackslash{}(\textbackslash{}lambda = 75.63 \textbackslash{}, \textbackslash{}text\{nm\}\textbackslash{}).
- From \textbackslash{}( n = 3 \textbackslash{}) to \textbackslash{}( n = 1 \textbackslash{}), \textbackslash{}(\textbackslash{}lambda = 79.76 \textbackslash{}, \textbackslash{}text\{nm\}\textbackslash{}).
- From \textbackslash{}( n = 2 \textbackslash{}) to \textbackslash{}( n = 1 \textbackslash{}), \textbackslash{}(\textbackslash{}lambda = 94.54 \textbackslash{}, \textbackslash{}text\{nm\}\textbackslash{}).

This diagram illustrates how electrons transition between energy states, releasing specific wavelengths of light as they move to the \textbackslash{}( n = 1 \textbackslash{}) state.

\#\# Dataset metadata

Quantum Phenomena; Multi-Formula Reasoning, Physical Model Grounding Reasoning

\paragraph{Recorded answer/context.}
D: D: \textbackslash{}( 378 \textbackslash{}text\{ nm\} \textbackslash{})

\paragraph{Figure/image path.}
/root/proposal\_for\_physic/science-mango/phyx\_mini\_run/phyx\_data/test\_image/519.png

\paragraph{Formalization target.}
create a compiling Lean file with sorry bodies at `PhyXMiniProblems/problem\_phyx\_mini\_0519.lean`. The Lean declarations must preserve the physical quantities, dimensions or dimensional roles, figure labels, governing-law hypotheses, and final relation expressed by this problem.
Use Mathlib/Physlib names found through LeanExplore where available. If a physics API is missing, introduce faithful local abstractions rather than scalar placeholder aliases.

\begin{theorem}[Physics formalization target]
\label{thm:physics:phyx_mini_0519:target}
\lean{PhyXMiniProblems.ProblemPhyXMini0519.emittedPhotonWavelengthN4ToN2_roundsTo378Nanometers}
\uses{def:physics:phyx-mini-0519:phyxminiproblems-problemphyxmini0519-oneelectronatomexperiment, def:physics:phyx-mini-0519:phyxminiproblems-problemphyxmini0519-matchesoneelectronatomscenario, def:physics:phyx-mini-0519:phyxminiproblems-problemphyxmini0519-matchessuppliedenergyleveldiagram, def:physics:phyx-mini-0519:phyxminiproblems-problemphyxmini0519-matchesionizationenergymeasurement, def:physics:phyx-mini-0519:phyxminiproblems-problemphyxmini0519-hasphysicalatomicspectrum, def:physics:phyx-mini-0519:phyxminiproblems-problemphyxmini0519-satisfiesatomicphotonlaws, def:physics:phyx-mini-0519:phyxminiproblems-problemphyxmini0519-roundstonearestwholenanometer}
The assigned autoformalize agent should translate this physics problem into Lean declarations in the covered file, with theorem and lemma proof bodies written as `by sorry`.
\end{theorem}
\begin{proof}
This is an autoformalization task, not a proof task. Produce statements that can later be proved by the physics prover without weakening the physical model.
\end{proof}

% --- Archon declaration topology begin ---
\section{Declaration topology}
\begin{definition}[Length Quantity]
  \label{def:physics:phyx-mini-0519:phyxminiproblems-problemphyxmini0519-lengthquantity}
  \lean{PhyXMiniProblems.ProblemPhyXMini0519.LengthQuantity}
  A unit-independent, nonnegative physical quantity with length dimension
\end{definition}
\begin{proof}
  This is the declared model definition of Length Quantity.
\end{proof}
\begin{definition}[length In Nanometers]
  \label{def:physics:phyx-mini-0519:phyxminiproblems-problemphyxmini0519-lengthinnanometers}
  \lean{PhyXMiniProblems.ProblemPhyXMini0519.lengthInNanometers}
  \uses{def:physics:phyx-mini-0519:phyxminiproblems-problemphyxmini0519-lengthquantity}
  Read a physical length as a real number of nanometres
\end{definition}
\begin{proof}
  This is the declared model definition of length In Nanometers.
\end{proof}
\begin{definition}[energy In Electron Volts]
  \label{def:physics:phyx-mini-0519:phyxminiproblems-problemphyxmini0519-energyinelectronvolts}
  \lean{PhyXMiniProblems.ProblemPhyXMini0519.energyInElectronVolts}
  Read a physical energy as a real number of electron volts. The quotient is dimensionless: both numerator and denominator are first represented in SI, and the denominator is Physlib's unit-independent one-electron-volt energy
\end{definition}
\begin{proof}
  This is the declared model definition of energy In Electron Volts.
\end{proof}
\begin{definition}[planck Times Light In Electron Volt Nanometers]
  \label{def:physics:phyx-mini-0519:phyxminiproblems-problemphyxmini0519-plancktimeslightinelectronvoltnanometers}
  \lean{PhyXMiniProblems.ProblemPhyXMini0519.planckTimesLightInElectronVoltNanometers}
  The standard product h c, read in electron-volt nanometres. Physlib supplies the reduced Planck constant in joule-seconds and the speed of light in metres per second; 2π converts ℏ to h, and 10\textasciicircum{}9 converts metres to nanometres
\end{definition}
\begin{proof}
  This is the declared model definition of planck Times Light In Electron Volt Nanometers.
\end{proof}
\begin{definition}[Atomic Level]
  \label{def:physics:phyx-mini-0519:phyxminiproblems-problemphyxmini0519-atomiclevel}
  \lean{PhyXMiniProblems.ProblemPhyXMini0519.AtomicLevel}
  The five principal-quantum-number labels printed on the supplied diagram
\end{definition}
\begin{proof}
  This is the declared model definition of Atomic Level.
\end{proof}
\begin{definition}[principal Quantum Number]
  \label{def:physics:phyx-mini-0519:phyxminiproblems-problemphyxmini0519-atomiclevel-principalquantumnumber}
  \lean{PhyXMiniProblems.ProblemPhyXMini0519.AtomicLevel.principalQuantumNumber}
  \uses{def:physics:phyx-mini-0519:phyxminiproblems-problemphyxmini0519-atomiclevel}
  The principal quantum number represented by an atomic-level label
\end{definition}
\begin{proof}
  This is the declared model definition of principal Quantum Number.
\end{proof}
\begin{definition}[Atomic Energy Level Diagram]
  \label{def:physics:phyx-mini-0519:phyxminiproblems-problemphyxmini0519-atomicenergyleveldiagram}
  \lean{PhyXMiniProblems.ProblemPhyXMini0519.AtomicEnergyLevelDiagram}
  \uses{def:physics:phyx-mini-0519:phyxminiproblems-problemphyxmini0519-lengthquantity, def:physics:phyx-mini-0519:phyxminiproblems-problemphyxmini0519-atomiclevel}
  The physically meaningful content read from an atomic energy-level diagram. The optional wavelength label is itself a physical length, rather than an untyped scalar attached to an arrow
\end{definition}
\begin{proof}
  This is the declared model definition of Atomic Energy Level Diagram.
\end{proof}
\begin{definition}[One Electron Atom Experiment]
  \label{def:physics:phyx-mini-0519:phyxminiproblems-problemphyxmini0519-oneelectronatomexperiment}
  \lean{PhyXMiniProblems.ProblemPhyXMini0519.OneElectronAtomExperiment}
  \uses{def:physics:phyx-mini-0519:phyxminiproblems-problemphyxmini0519-lengthquantity, def:physics:phyx-mini-0519:phyxminiproblems-problemphyxmini0519-atomiclevel, def:physics:phyx-mini-0519:phyxminiproblems-problemphyxmini0519-atomicenergyleveldiagram}
  Independent quantities of the hypothetical atom and its emission experiment. The n = 4 to n = 2 wavelength is an observable field; it is not defined to be the recorded answer. The laws below constrain all downward transitions
\end{definition}
\begin{proof}
  This is the declared model definition of One Electron Atom Experiment.
\end{proof}
\begin{definition}[Matches One Electron Atom Scenario]
  \label{def:physics:phyx-mini-0519:phyxminiproblems-problemphyxmini0519-matchesoneelectronatomscenario}
  \lean{PhyXMiniProblems.ProblemPhyXMini0519.MatchesOneElectronAtomScenario}
  \uses{def:physics:phyx-mini-0519:phyxminiproblems-problemphyxmini0519-energyinelectronvolts, def:physics:phyx-mini-0519:phyxminiproblems-problemphyxmini0519-atomiclevel, def:physics:phyx-mini-0519:phyxminiproblems-problemphyxmini0519-oneelectronatomexperiment}
  The qualitative one-electron, ground-state, and bound-level scenario
\end{definition}
\begin{proof}
  This is the declared model definition of Matches One Electron Atom Scenario.
\end{proof}
\begin{definition}[Matches Supplied Energy Level Diagram]
  \label{def:physics:phyx-mini-0519:phyxminiproblems-problemphyxmini0519-matchessuppliedenergyleveldiagram}
  \lean{PhyXMiniProblems.ProblemPhyXMini0519.MatchesSuppliedEnergyLevelDiagram}
  \uses{def:physics:phyx-mini-0519:phyxminiproblems-problemphyxmini0519-lengthinnanometers, def:physics:phyx-mini-0519:phyxminiproblems-problemphyxmini0519-atomiclevel, def:physics:phyx-mini-0519:phyxminiproblems-problemphyxmini0519-oneelectronatomexperiment}
  Primary-image readout. It records all five horizontal level lines, the four downward arrows ending at n = 1, and their wavelength labels. In particular, the requested n = 4 to n = 2 arrow and wavelength are not shown in the figure and therefore do not occur here as measured answer data
\end{definition}
\begin{proof}
  This is the declared model definition of Matches Supplied Energy Level Diagram.
\end{proof}
\begin{definition}[Matches Ionization Energy Measurement]
  \label{def:physics:phyx-mini-0519:phyxminiproblems-problemphyxmini0519-matchesionizationenergymeasurement}
  \lean{PhyXMiniProblems.ProblemPhyXMini0519.MatchesIonizationEnergyMeasurement}
  \uses{def:physics:phyx-mini-0519:phyxminiproblems-problemphyxmini0519-energyinelectronvolts, def:physics:phyx-mini-0519:phyxminiproblems-problemphyxmini0519-oneelectronatomexperiment}
  The independently observed 17.50 eV ground-state ionization energy
\end{definition}
\begin{proof}
  This is the declared model definition of Matches Ionization Energy Measurement.
\end{proof}
\begin{definition}[Has Physical Atomic Spectrum]
  \label{def:physics:phyx-mini-0519:phyxminiproblems-problemphyxmini0519-hasphysicalatomicspectrum}
  \lean{PhyXMiniProblems.ProblemPhyXMini0519.HasPhysicalAtomicSpectrum}
  \uses{def:physics:phyx-mini-0519:phyxminiproblems-problemphyxmini0519-lengthinnanometers, def:physics:phyx-mini-0519:phyxminiproblems-problemphyxmini0519-energyinelectronvolts, def:physics:phyx-mini-0519:phyxminiproblems-problemphyxmini0519-atomiclevel, def:physics:phyx-mini-0519:phyxminiproblems-problemphyxmini0519-atomiclevel-principalquantumnumber, def:physics:phyx-mini-0519:phyxminiproblems-problemphyxmini0519-oneelectronatomexperiment}
  Strict ordering and positivity conditions for the physical spectrum
\end{definition}
\begin{proof}
  This is the declared model definition of Has Physical Atomic Spectrum.
\end{proof}
\begin{definition}[Satisfies Atomic Photon Laws]
  \label{def:physics:phyx-mini-0519:phyxminiproblems-problemphyxmini0519-satisfiesatomicphotonlaws}
  \lean{PhyXMiniProblems.ProblemPhyXMini0519.SatisfiesAtomicPhotonLaws}
  \uses{def:physics:phyx-mini-0519:phyxminiproblems-problemphyxmini0519-lengthinnanometers, def:physics:phyx-mini-0519:phyxminiproblems-problemphyxmini0519-energyinelectronvolts, def:physics:phyx-mini-0519:phyxminiproblems-problemphyxmini0519-plancktimeslightinelectronvoltnanometers, def:physics:phyx-mini-0519:phyxminiproblems-problemphyxmini0519-atomiclevel, def:physics:phyx-mini-0519:phyxminiproblems-problemphyxmini0519-atomiclevel-principalquantumnumber, def:physics:phyx-mini-0519:phyxminiproblems-problemphyxmini0519-oneelectronatomexperiment}
  The two governing relations used by the calculation: ionization energy is the gap from the n = 1 level to the continuum; every downward emitted photon obeys ΔE λ = h c, here expressed in electron volts and nanometres. The transition law is uniform in the two levels and contains no special value for the requested n = 4 to n = 2 wavelength
\end{definition}
\begin{proof}
  This is the declared model definition of Satisfies Atomic Photon Laws.
\end{proof}
\begin{definition}[Rounds To Nearest Whole Nanometer]
  \label{def:physics:phyx-mini-0519:phyxminiproblems-problemphyxmini0519-roundstonearestwholenanometer}
  \lean{PhyXMiniProblems.ProblemPhyXMini0519.RoundsToNearestWholeNanometer}
  \uses{def:physics:phyx-mini-0519:phyxminiproblems-problemphyxmini0519-lengthquantity, def:physics:phyx-mini-0519:phyxminiproblems-problemphyxmini0519-lengthinnanometers}
  A physical wavelength rounds to the displayed whole number of nanometres when its nanometre readout lies in the corresponding half-open unit interval
\end{definition}
\begin{proof}
  This is the declared model definition of Rounds To Nearest Whole Nanometer.
\end{proof}
% --- Archon declaration topology end ---
% --- Archon physics formalization source end ---
